\section{The Small-Correlation Regime}
\label{app:low}

This appendix proves CK for \(0\le\rho\le\rho_*\), at every bias,
where \(\rho_*=\tanh(31/20)\) is the common boundary.
Its mechanism is an energy gap. A rule with one large singleton
coefficient is covered by the local lemma. Otherwise its total singleton
energy is bounded away from one. Two quadratics tangent to binary
entropy convert that energy gap into the required information bound.
This is the local-optimality/Fourier-gap strategy developed in
Samorodnitsky's and Yu's work~\cite{Sam,YuPhi,Yu}. The distinction in
this appendix is its self-contained quantitative input: the cap below
is proved by sections and sign sums, and the entropy comparison keeps
the output mean. Yu's published balanced range is not being used as an
all-bias black box.

Write \(r=\rho\) in this appendix. For a cube function \(v\), its
\emph{singleton energy} is
\(W_1(v)=\sum_i\widehat v(\{i\})^2\). Let
\(a=\max_i|\widehat f(\{i\})|\) be the largest singleton coefficient
in absolute value. Use the cutoff \(a_0=5/8\) and the energy cap
\[
                   M=1-a_0+a_0^2=\left(\frac78\right)^2.
\]
These choices match two bounds: splitting on the strongest coordinate
gives \(1-a_0+a_0^2\), while the sign-sum estimate below gives
\(7/8\). The cutoff makes their energy bounds agree exactly.
We shall start the entropy comparison at \(r_0=5/8\) and end at
\(r_1=\rho_*\).
Contraction will then cover all correlations below \(r_0\).

\subsection{An elementary bound on singleton energy}
\label{app:fourier-cap-proof}

We prove the Fourier estimate needed below, including the estimate for
sums of independent signs on which it rests. Let \(f\) be any function
on the cube with values in \(\{-1,1\}\), and put
\[
 a_i=\widehat f(\{i\}),\qquad a=\max_i|a_i|,\qquad
 W=\sum_i a_i^2=W_1(f).
\]
Thus \(a\) measures the strongest correlation with a single input bit,
and \(W\) is the total squared correlation with all individual bits.
We will show
\begin{equation}
                 a\le\frac58
                 \quad\Longrightarrow\quad W\le 1-a_0+a_0^2=M.
 \label{eq:elementary-fourier-cap}
\end{equation}
The estimate holds at every mean. It separates into two simple cases:
one coefficient is moderately large, or every coefficient is small.
This is a quantitative instance of the level-one rigidity underlying
FKN~\cite{FKN}, tailored to a fixed cutoff rather than an asymptotic
closeness theorem. For the broader sign-sum problem see~\cite{KSTJ},
and for a sharp small-coefficient asymptotic see~\cite{YuSharp}.
The particular constants needed here are established below.

\subsubsection{A moderately large coefficient leaves little other energy}

Orient a chosen input bit so that its coefficient is \(a\ge0\), and
let \(f_+,f_-\) be the two sections obtained by fixing that bit.
On the remaining cube define
\[
 s=\frac{f_++f_-}{2},\qquad d=\frac{f_+-f_-}{2}.
\]
Pointwise, \(s^2+d^2=1\) and \(d\in\{-1,0,1\}\). Also
\(\E d=a\), so \(\E d^2\ge a\). The other singleton coefficients
are the singleton coefficients of \(s\). Parseval therefore gives
\begin{equation}
 W=a^2+W_1(s)\le a^2+\E s^2\le a^2+1-a.
 \label{eq:elementary-singleton-section}
\end{equation}
For \(3/8\le a\le5/8\), the convex quadratic on the right is at
most its equal endpoint values:
\[
 \left(\frac38\right)^2+1-\frac38
 =\left(\frac58\right)^2+1-\frac58
 =\left(\frac78\right)^2=M.
\]
It remains to handle \(a\le3/8\).

\subsubsection{Spread-out coefficients cannot have absolute mean near one}

We use the following elementary moment bound. If
\(Z=\sum_i b_iX_i\), where the \(X_i\) are independent uniform signs, then
\begin{equation}
 \sum_i b_i^2=1,\qquad \max_i b_i^2\le\frac15
 \quad\Longrightarrow\quad \E|Z|<\frac78.
 \label{eq:elementary-sign-sum}
\end{equation}
The normalization gives \(\E Z^2=1\), so Cauchy--Schwarz alone
would give only \(\E|Z|\le1\). The condition on the largest
summand gives a strict gap. We prove it with one polynomial above
absolute value; no Gaussian approximation or tail estimate is needed.

\paragraph{Contact polynomial.}
Let \(p\) be the polynomial of degree at most five that matches
\(\sqrt y\) and its derivative at the three points
\[
                         y=9/16,\quad 4,\quad 9.
\]
These are convenient contacts for one fixed bound, not parameters of
the channel. The contact conditions determine \(p\) uniquely:
the difference of two such polynomials would have six zeros, counting
multiplicity, despite having degree at most five.
For complete arithmetic, put \(D_{\rm mom}=4042912500\). The polynomial is
\begin{equation}
 \begin{split}
 D_{\rm mom}p(y)={}&637696y^5-17928800y^4\\
                  &+191446315y^3-1008169920y^2\\
                  &+3660186285y+1260006516.
 \end{split}
 \label{eq:moment-majorant}
\end{equation}
Substitution verifies the six contact conditions.

To prove that the contact polynomial is a global majorant, fix \(y>0\)
away from the contact points. Subtract from \(\sqrt s-p(s)\) the multiple
of \((s-9/16)^2(s-4)^2(s-9)^2\) that makes it zero at \(s=y\).
There are seven zeros counting multiplicity. Applying Rolle's theorem
six times gives, for some \(\xi>0\),
\[
 \sqrt y-p(y)=\frac{(\sqrt{\,\cdot\,})^{(6)}(\xi)}{6!}
                (y-9/16)^2(y-4)^2(y-9)^2\le0,
\]
since \((\sqrt{\,\cdot\,})^{(6)}(\xi)=-945/(64\xi^{11/2})<0\).
The contact points and zero follow by continuity. Thus
\[
                              |Z|\le p(Z^2).
\]
This sign argument is the reason for using a contact polynomial.

\paragraph{Moment evaluation.}
Write \(\sigma_j=\sum_i b_i^j\) for \(j=4,6,8,10\).
Expansion of the sum, keeping only even multiplicities, gives
\begin{align*}
 \E Z^2&=1, & \E Z^4&=3-2\sigma_4,\\
 \E Z^6&=15-30\sigma_4+16\sigma_6,\\
 \E Z^8&=105-420\sigma_4+140\sigma_4^2\\
        &\quad+448\sigma_6-272\sigma_8,\\
 \E Z^{10}&=945-6300\sigma_4+6300\sigma_4^2\\
 &\quad+10080\sigma_6-6720\sigma_4\sigma_6\\
 &\quad-12240\sigma_8+7936\sigma_{10}.
\end{align*}
One systematic way to obtain these identities is to differentiate
\(\E e^{tZ}=\prod_i\cosh(b_it)\) at zero. The coefficients needed are
\begin{align*}
 \log\cosh t={}&\frac{t^2}{2}-\frac{t^4}{12}+\frac{t^6}{45}\\
 &-\frac{17t^8}{2520}+\frac{31t^{10}}{14175}+O(t^{12});
\end{align*}
they follow by integrating the Taylor series of \(\tanh t\).

Put \(t=\sigma_4\) in the rest of this paragraph. The largest-weight
condition implies
\begin{align*}
 0\le t&\le\frac15,& \sigma_6&\le\frac t5,\\
 \sigma_{10}&\le\frac{\sigma_8}{5},& \sigma_8&\ge t^3.
\end{align*}
For the last inequality, the numbers \(b_i^2\) are probability
weights: Jensen for the convex map \(x\mapsto x^3\) gives
\(\sum_i b_i^2(b_i^2)^3\ge(\sum_i b_i^2b_i^2)^3\).
Substitute the moment identities into \eqref{eq:moment-majorant}, then
apply these four inequalities. The result is the single cubic bound
\[
 D_{\rm mom}\E p(Z^2)\le T_{\rm mom}(t),
\]
where
\begin{align*}
 T_{\rm mom}(t)={}&3487476486+77364454t\\
 &+650389376t^2-\frac{9583071744}{5}t^3.
\end{align*}
To check the substitution directions, the coefficient of \(\sigma_6\)
before using \(\sigma_6\le t/5\) is
\(1459014320-4285317120t>0\) on \([0,1/5]\).
After \(\sigma_{10}\le\sigma_8/5\), the coefficient of
\(\sigma_8\) is \(-9583071744/5<0\), so Jensen is used in the correct
direction.

Only one value of this cubic needs checking. Its derivative is concave,
and is positive at both endpoints of \([0,1/5]\):
\begin{align*}
 T_{\rm mom}'(0)&=77364454>0,\\
 T_{\rm mom}'(1/5)&=\frac{13440810318}{125}>0.
\end{align*}
Hence \(T_{\rm mom}(t)\le T_{\rm mom}(1/5)\). The remaining exact comparison is
\[
 7D_{\rm mom}-8T_{\rm mom}(1/5)=\frac{119582002252}{625}>0.
\]
It follows that \(\E|Z|\le\E p(Z^2)<7/8\), proving
\eqref{eq:elementary-sign-sum}.

\subsubsection{Return to the Boolean function}

Suppose \(a\le3/8\) and \(W\ge M=(7/8)^2\), and normalize
\(b_i=a_i/\sqrt W\). Then
\[
 \sum_i b_i^2=1,\qquad
 \max_i b_i^2\le\frac{(3/8)^2}{(7/8)^2}=\frac9{49}<\frac15.
\]
Booleanity supplies the bridge to the moment bound:
\[
 \sqrt W=\E\!\left[f\sum_i b_iX_i\right]
       \le\E\left|\sum_i b_iX_i\right|<\frac78.
\]
This contradicts \(\sqrt W\ge\sqrt M=7/8\) directly.
Thus small coefficients give \(W<M\). Together with the section
bound on \([3/8,5/8]\), this proves
\eqref{eq:elementary-fourier-cap}.

\subsection{The same cutoff covers larger coefficients}
\label{app:low-local-check}

Recall the local margin \(M_{\rm loc}\) from
Section~\ref{app:local-propagation}. At our common boundary
\(r_1=\tanh(31/20)\), the fixed comparison in
Section~\ref{app:low-arithmetic} gives
\[
 M_{\rm loc}(r_1,5/8)
  =(1-r_1^2)h(5r_1/8)-(1-5r_1^2/8)h(r_1)>0.
\]
The propagation argument there extends this one check to every
\(r\le r_1\) and \(a\ge5/8\). Thus every rule is covered:
a coefficient at least \(5/8\) invokes the local lemma; otherwise
the singleton energy is at most \((7/8)^2\).


\subsection{Mean-Sensitive Energy Bound}

The Fourier estimate just proved says
\begin{equation}
 a\le a_0\quad\Longrightarrow\quad W_1(f)\le M.
 \label{eq:fourier-cap}
\end{equation}
The local check in Section~\ref{app:low-local-check}, followed by
Section~\ref{app:local-propagation}, handles \(a\ge a_0\) throughout
\(r\le r_1\).

For a biased rule, introduce one new sign bit and form
\[
 G(1,x)=f(x),\qquad G(-1,x)=-f(-x).
\]
The new rule is balanced. Its singleton coefficient on the new bit is
\(m=\E f\); on the old bits it has the original singleton coefficients.
Consequently, if \(|m|\le a_0\) and \(a\le a_0\),
\begin{equation}
                         m^2+W_1(f)=W_1(G)\le M.
 \label{eq:lift}
\end{equation}
Only this Fourier relation is needed from \(G\).
Larger biases are covered directly by entropy contraction:
\[
 I_f(r)\le r^2h(m)<r^2/2\le\phi(r),
\]
Indeed, the first two terms of the entropy series give, when
\(|m|\ge a_0=5/8\),
\begin{align*}
 h(m)&\le h(a_0)\le L-a_0^2/2-a_0^4/12\\
 &<7/10-a_0^2/2-a_0^4/12<1/2.
\end{align*}
The point \(r=0\) is immediate.

\subsection{Tangent Quadratic Minorants}

\begin{remark}[Role of the quadratic minorant]
The first two moments of the noisy rule have simple Fourier formulas.
A quadratic below entropy lets us use those formulas directly. We
construct the quadratic by making it touch entropy, rather than by
fitting its coefficients to a numerical grid.
\end{remark}

For a contact point \(0<t<1\) and an entropy argument \(0\le s\le1\), define
\begin{align*}
 A_t&=\frac{\log(2/(1+t))}{(1-t)^2},\\
 z_t&=2t-\frac{g(t)}{A_t},\\
 Q_t(s)&=A_t(1-s)(1+s-z_t).
\end{align*}
The coefficient \(A_t\) sets the scale; \(z_t\) sets the shape.
They are determined by contact at \(t\) and vanishing at one. The identity
\(h(t)-(1-t)g(t)=\log(2/(1+t))\) gives
\[
 Q_t(t)=h(t),\qquad Q_t'(t)=h'(t),\qquad Q_t(1)=h(1)=0.
\]
Moreover, every such quadratic is a global lower bound:
\begin{equation}
                         Q_t(s)\le h(s)\qquad(0\le s\le1).
 \label{eq:entropy-tangent-family}
\end{equation}
To see this, put \(D_t=h-Q_t\). Its second derivative
\(D_t''(s)=2A_t-1/(1-s^2)\) strictly decreases. It must be positive
at \(t\); otherwise \(D_t\) would decrease strictly from its zero at
\(t\) to its zero at one. Thus \(D_t\) decreases to zero before
\(t\), then rises and falls back to zero on \([t,1]\).
It stays nonnegative throughout. This proves the entire family at once.

Use the two contact points \(t_0=2/5\) and \(t_1=2/3\).
For \(j=0,1\), abbreviate \(A_j=A_{t_j}\) and \(z_j=z_{t_j}\).
Define \(\alpha,\beta\) by the line
through the two points \((z_j,1/A_j)\):
\begin{align*}
 \beta&=\frac{1/A_0-1/A_1}{z_1-z_0},\\
 \alpha&=\frac1{A_0}+\beta z_0,\\
 c_{\rm ent}(r)&=\frac1{\alpha-\beta r^2}.
\end{align*}
Thus the coefficient is determined by the two entropy tangencies.
The fixed comparisons at the end of this appendix give
\(z_0<r_0^2\) and \(r_1^2<z_1\).
For any \(r\in[r_0,r_1]\), write \(r^2=\mu z_0+(1-\mu)z_1\),
where \(0\le\mu\le1\). The weights
\(\mu c_{\rm ent}/A_0\) and \((1-\mu)c_{\rm ent}/A_1\) sum to one.
Their convex combination of \eqref{eq:entropy-tangent-family} gives
\begin{equation}
 h(s)\ge c_{\rm ent}(r)(1-s)(1+s-r^2).
 \label{eq:minorant}
\end{equation}
This is the only entropy minorant used below.

\subsection{From Entropy to Fourier Energy}

For each Fourier degree \(k\), put
\(W_k=\sum_{|S|=k}\widehat f(S)^2\). Since \(f\) takes only the
values \(-1,1\),
\[
 \E|T_r f|\ge\E[fT_r f]=\sum_{k=0}^n r^kW_k.
\]
Apply \eqref{eq:minorant} to \(|T_r f|\), and use Parseval for its
second moment. All higher-degree contributions have the favorable sign:
\(r^{k+2}-r^{2k}\ge0\) for \(k\ge2\). Hence
\[
 \E h(T_r f)\ge
 c_{\rm ent}(r)\bigl[1-r^2-(1-r^2)m^2-r^2(1-r)W_1\bigr].
\]
Insert \eqref{eq:lift} and \(\phi(m)\ge m^2/2\):
\begin{align}
 \E h(T_r f)+\phi(m)\ge{}&
 c_{\rm ent}(r)(1-r)(1+r-Mr^2)\notag\\
 &+\bigl[1/2-c_{\rm ent}(r)(1-2r^2+r^3)\bigr]m^2.
 \label{eq:low-credit}
\end{align}
The mean coefficient is positive. After multiplication by
\(2/c_{\rm ent}(r)>0\), it becomes
\[
 \alpha-2+(4-\beta)r^2-2r^3.
\]
This increases on \([r_0,1]\), since \(\beta<1\). At \(r_0=5/8\),
the bounds \(\alpha>53/40\) and \(\beta<17/20\) give a value
larger than \(43/640>0\). We may therefore discard the mean term.

\subsection{One change of curvature controls the remaining gap}

The information inequality now reduces to the positivity of
\[
 B_0(r)=(1-r)(1+r-Mr^2)-(\alpha-\beta r^2)h(r).
\]
We prove that \(B_0\) decreases, so a single comparison at \(r_1\)
finishes the interval. The reason is a simple derivative shape:
\begin{equation}
 B_0''''(r)=
 \frac{2[\alpha(1+3r^2)-\beta(r^4-3r^2+6)]}{(1-r^2)^3}.
 \label{eq:low-curvature-shape}
\end{equation}
The numerator increases as a function of \(r^2\). Thus \(B_0''''\)
can change sign only from negative to positive. Since
\(B_0'''(r_0)<0\), the same is true of \(B_0'''\).
It follows that \(B_0''\) has no interior maximum on any subinterval.

The fixed comparisons below establish
\[
 \begin{gathered}
 B_0'(r_0)<-1/16,\qquad B_0''(r_0)<1/3,\qquad B_0'''(r_0)<0,\\
 B_0''(4/5)<0,\qquad B_0''(r_1)<0.
 \end{gathered}
\]
The endpoint property gives \(B_0''<1/3\) up to \(4/5\), and
\(B_0''<0\) thereafter. Consequently
\begin{align*}
 B_0'(r)&<-\frac1{16}+\frac{4/5-5/8}{3}
          =-\frac1{240}<0,\\
 &\hspace{8em}r_0\le r\le4/5,
\end{align*}
and \(B_0'\) decreases after \(4/5\).
Therefore \(B_0(r)\ge B_0(r_1)>0\) throughout the interval.
Equation~\eqref{eq:low-credit} now gives CK on \([r_0,r_1]\).

\subsection{Contract from the starting point to zero}

The same estimate leaves extra slack at \(r_0\):
\[
 c_{\rm ent}(r_0)(1-r_0)(1+r_0-Mr_0^2)+r_0^2/2-L>0.
\]
Thus \(I_f(r_0)<r_0^2/2\). Entropy contraction gives
\[
 I_f(r)\le(r/r_0)^2I_f(r_0)\le r^2/2\le\phi(r)
                   \qquad(0\le r\le r_0).
\]
The strict starting slack matters: contracting a bound equal to
\(\phi(r_0)\) would not suffice, since \(\phi(r)/r^2\) increases.
The local and large-bias alternatives were already covered.
This completes the all-bias proof for \(0\le r\le\rho_*\).

\subsection{The fixed comparisons}
\label{app:low-arithmetic}

All coefficients above were defined by tangency. The following bounds
only check finitely many values; they do not specify a fitted entropy
curve. For \(1\le y\le2\), put \(z=(y-1)/(y+1)\) and define
\begin{align*}
 \ell_-(y)&=2\sum_{j=0}^{11}\frac{z^{2j+1}}{2j+1},\\
 \ell_+(y)&=\ell_-(y)+\frac{2z^{25}}{25(1-z^2)}.
\end{align*}
The logarithm series gives
\[
 \ell_-(y)\le\log y\le\ell_+(y).
\]
For other positive rational arguments, first extract a power of two.
Only the following logarithms occur:
\[
 \log(10/7),\qquad \log(6/5),\qquad \log(7/3),\qquad \log5.
\]
Substituting the displayed sums gives
\[
 \alpha>\frac{53}{40},\qquad 0<\beta<\frac{17}{20},\qquad
 z_0<r_0^2,\quad r_1^2<z_1.
\]
These are the bounds used for the mean term, derivative shape, and
interpolation domain. In the fixed checks below, evaluate \(\alpha\)
and \(\beta\) directly from their defining logarithms and the displayed
finite sums; no rounded coefficient values are needed.
For scale only, the same substitutions with outward rounding give
\begin{align*}
 1.3258315440444&<\alpha<1.3258315440447,\\
 0.8499134388946&<\beta<0.8499134388950.
\end{align*}
These decimal enclosures help a reader reproduce the calculation, but
the proof uses the finite rational sums above.

For completeness, the derivatives needed in the five fixed checks are
\begin{align*}
 B_0'(r)={}&-2(1+M)r+3Mr^2\\
 &+(\alpha-\beta r^2)g(r)+2\beta r h(r),\\
 B_0''(r)={}&-2(1+M)+6Mr+\frac{\alpha-\beta r^2}{1-r^2}\\
 &-4\beta r g(r)+2\beta h(r),\\
 B_0'''(r)={}&6M-6\beta g(r)-\frac{6\beta r}{1-r^2}\\
 &+\frac{2r(\alpha-\beta r^2)}{(1-r^2)^2}.
\end{align*}
At rational arguments, use \(g(r)=\tfrac12\log((1+r)/(1-r))\)
and the definition of \(h(r)\), replacing each logarithm by
\(\ell_-\) or \(\ell_+\) in the unfavorable direction.
At the common boundary, the same exponent gives all three values:
\begin{align*}
 r_1&=\frac{e^{31/10}-1}{e^{31/10}+1},&
 g(r_1)&=\frac{31}{20},\\
 h(r_1)&=\log(1+e^{-31/10})+\frac{31/10}{1+e^{31/10}}.
\end{align*}
The exponential sum and remainder in Section~\ref{app:fixed-arithmetic}
bound the exponential. Substitution in these identities, followed by
the logarithm bounds, gives the endpoint comparisons and
\(r_1^2<z_1\). All substitutions are in the unfavorable direction
in the displayed formula for \(B_0\) and its derivatives. For the local
margin, use the same bounds on \(r_1\) and the logarithm formula for
\(h(5r_1/8)\). These finite sums establish the table, with
\(M=(7/8)^2\):
\begin{center}\small
\begin{tabular}{@{}p{.71\linewidth}r@{}}
\toprule
Fixed expression & Bound\\ \midrule
\(M_{\rm loc}(r_1,5/8)\) & \(>1/1400\)\\
\(B_0'(5/8)\) & \(<-17/250\)\\
\(B_0''(5/8)\) & \(<7/30\)\\
\(B_0'''(5/8)\) & \(<-1\)\\
\(B_0''(4/5)\) & \(<-1/9\)\\
\(B_0''(r_1)\) & \(<-1/10\)\\
\(B_0(r_1)\) & \(>1/4000\)\\
Starting contraction slack at \(r_0=5/8\) & \(>1/500\)\\
\bottomrule
\end{tabular}
\end{center}
These stronger rational margins are obtained by the same unfavorable
finite-sum substitutions and exceed the weaker signs used in the proof.
Every check uses a fixed logarithm or exponential value. The interval
argument is the derivative shape in \eqref{eq:low-curvature-shape}, not
a grid of arithmetic checks.
